% !TeX root = ../main.tex
% !TeX encoding = UTF-8
% !TeX spellcheck = none

% Reference identities and placements retained from the verified Chinese manuscript.
\section{Related Work}\label{sec:related-work}

\subsection{Spatiotemporal Representations and Temporal Pattern Disentanglement}\label{sec:frequency-related}

Traffic flow forecasting combines spatial dependency modeling with temporal encoding. Graph neural networks make this structure explicit by modeling associations between traffic nodes~\cite{ref29}. DCRNN represents propagation as diffusion on a directed graph and combines diffusion convolution with recurrent units~\cite{ref01}; STGCN couples graph and temporal convolution~\cite{ref02}; and Graph WaveNet learns latent relations with an adaptive adjacency matrix and dilated convolutions~\cite{ref03}. STID adds spatial and temporal identity embeddings to distinguish nodes and time slots~\cite{ref28}. These models define the standard spatiotemporal representation used by regional forecasters.

Attention and adaptive relation learning further refine this representation. ASTGCN and GMAN use spatial and temporal attention to focus on relevant locations and times~\cite{ref04,ref05}. MTGNN and AGCRN learn node relations from data~\cite{ref06,ref07}, while STSGCN and STFGNN construct local spatiotemporal graphs from node--time associations and sequence similarity~\cite{ref08,ref09}. These models improve dependency learning while maintaining a unified representation of the traffic history.

Temporal pattern disentanglement organizes encoding around frequency or periodic structure. STWave uses a discrete wavelet transform to separate low- and high-frequency components~\cite{ref11}; TEDDCF separates trend and event components and assigns them different spatial and temporal paths~\cite{ref18}; and other methods distinguish periodic from aperiodic variation~\cite{ref37}. HyperD represents periodic patterns with daily and weekly embeddings, models residual fluctuations in the frequency domain, and aligns the two views~\cite{ref22}. These approaches make the base representation more responsive to dynamics at different scales.

These methods strengthen the base representation. They leave a complementary question: how should information from different sources correct the prediction when traffic departs from routine behavior?

\subsection{Heterogeneous Evidence and Expert Modeling}\label{sec:experts-related}

Traffic dynamics vary across regions and time periods. ST-SSL uses spatial and temporal self-supervised tasks to strengthen representations of this heterogeneity~\cite{ref30}. ST-ResNet uses three residual branches for temporal closeness, periodicity, and trends, and learns to fuse their outputs~\cite{ref14}.

Other methods explicitly distinguish traffic dynamics through disentanglement or expert mechanisms. D2STGNN separates diffusion and inherent signals through gating and residual decomposition and models the two types of dynamics separately~\cite{ref25}. TESTAM uses temporal, static-graph, and dynamic-graph experts, with gated Top-1 routing to adapt to changing spatiotemporal dependencies~\cite{ref26}.

These models organize heterogeneous dynamics with temporal branches, signal decomposition, or expert routing. Our design assigns each evidence source a correction role and gates its contribution to a shared base prediction using the current state.

Under deviations from routine patterns, the relevant distinction is the role of each evidence source. Directed time-lagged associations encode cross-region relations, whereas node self-evolution preserves local continuity; their different temporal meanings call for separate processing paths.

\subsection{Dynamic Dependencies and Historical Message Selection}\label{sec:messages-related}

Traffic forecasting also needs to select cross-region information as relations change with time. GMAN uses spatial and temporal attention and connects historical representations to future steps with transform attention~\cite{ref05}. PDFormer combines spatial self-attention with short- and long-range masks and uses delay-aware transformations to address temporal misalignment~\cite{ref10}. TrGNN extracts transition patterns from vehicle trajectories and examines nonrecurrent flow under peak-hour and exceptional-event conditions~\cite{ref40}. These studies show that traffic relations depend on temporal shifts and propagation delays as well as spatial distance~\cite{ref10,ref39}.

Adaptive graph learning infers latent spatial relations from data. MTGNN and AGCRN construct adjacency relations from learnable node embeddings~\cite{ref06,ref07}. These relations identify candidate interactions; a forecasting model must still determine whether a particular historical message is useful for the current context and forecast horizon.

Memory mechanisms provide another way to adapt spatial relations and retrieve historical patterns. MegaCRN generates graph structures from learnable meta-node memory~\cite{ref31}; RAST retrieves fine-grained patterns from an external memory bank and fuses them through cross-attention~\cite{ref23}. Cross-region forecasting therefore involves two linked decisions: which regions and lags form the candidate set, and which messages should be selected for the current target.

The usefulness of cross-region historical information depends on its context~\cite{ref46}. A neighboring region can carry different predictive value at different lags, temporal scales, and forecast steps. Our BCMS module addresses this selection problem by conditioning message scores on source and target backgrounds and on the forecast step.
